% ---------------------------------------------------------------------------
% content.tex — the preview deck (theme-independent)
%
% Every colour and every piece of chrome comes from the theme; nothing here
% names a colour or a font. This is the file the real slides will grow from.
%
%   % 1  title slide          — theme's title page template
%   % 2  layout preview       — text, bullets, table, callout
%   % 3  layout preview       — boxes, code, diagram
%   % 4  closing slide        — \closingSlide, overridable per theme
% ---------------------------------------------------------------------------

% ===========================================================================
% 1. Title slide
% ===========================================================================

\begin{frame}[plain]
  \titlepage
\end{frame}

% ===========================================================================
% 2. Layout preview — text, table, code
%
% Every element the deck can use, so a theme can be judged at a glance.
% Delete both preview slides once the section format is settled.
% ===========================================================================

\begin{frame}{Layout preview — text, table, code}
  \begin{columns}[t, onlytextwidth]
    \begin{column}{0.50\textwidth}
      {\color{rtmMuted}\scriptsize\bfseries SECTION HEADING}\par
      \vskip1mm
      \begin{itemize}
        \item First-level bullet with \concept{inline emphasis}
        \item A second bullet, and a \num{42.7\%} figure in accent colour
        \begin{itemize}
          \item A sub-item, quieter than its parent
        \end{itemize}
        \item Third bullet, closing the list
      \end{itemize}
      \vskip2mm
      {\color{rtmMuted}\scriptsize\bfseries EMPHASIS}\par
      \vskip1mm
      Body text in the normal ink colour, sized for projection, and
      \alert{alerted text} for something that must not be missed.
    \end{column}

    \begin{column}{0.46\textwidth}
      \begin{table}
        \centering
        \begin{tabular}{@{}lr@{}}
          \toprule
          \rtmtablehead{Quantity} & \rtmtablehead{Value} \\
          \midrule
          First row   & 1 \\
          Second row  & 2 \\
          Third row   & \num{3} \\
          \bottomrule
        \end{tabular}
      \end{table}
      \vskip1mm
      \begin{highlight}
        \textbf{Highlight box} — accent rule, for a point worth remembering.
      \end{highlight}
    \end{column}
  \end{columns}
  \themeSignature
\end{frame}

% ---------------------------------------------------------------------------
\begin{frame}[fragile]{Layout preview — boxes, code, diagram}
  \begin{columns}[t, onlytextwidth]
    \begin{column}{0.48\textwidth}
      \begin{tech}
        \textbf{Technical box} — primary rule, for an implementation note.
      \end{tech}
      \vskip2mm
      \begin{warn}
        \textbf{Warning box} — alert rule, for a caveat or a risk.
      \end{warn}
      \vskip3mm
      \key{Pill: one line, the slide's takeaway}
    \end{column}

    \begin{column}{0.50\textwidth}
\begin{lstlisting}[numbers=left]
/* Code block --- monospaced, */
int example(int n)   /* numbers optional */
{
    return n * n;
}
\end{lstlisting}
      \vskip1mm
      \begin{center}
      \begin{tikzpicture}[node distance=6mm and 5mm]
        \node[accent block] (a) {Start};
        \node[block, right=of a] (b) {Step};
        \node[muted block, right=of b] (c) {End};
        \draw[flow] (a) -- (b);
        \draw[flow dashed] (b) -- (c);
      \end{tikzpicture}
      \end{center}
    \end{column}
  \end{columns}
  \themeSignature
\end{frame}

% ===========================================================================
% 3. Closing slide
% ===========================================================================

\closingSlide

\end{document}
